% Chapter — Event detection. Follows the mandatory FR-29 six-section
% template: purpose; assumptions and domain of validity; derivation;
% discretization and implementation; validation evidence; references. The
% equation labels eq:events:screen, eq:events:ghat, eq:events:timer, and
% eq:events:apsis are echoed verbatim in comments in cpp/src/events.cpp and
% cpp/include/star/events.hpp (FR-29 equation-label-to-code traceability);
% renaming them breaks that trace.
\chapter{Event Detection}
\label{ch:events}

\section{Purpose}
\label{sec:events:purpose}

This chapter documents the event-detection framework (requirement FR-12):
scalar event functions $g(t, \vv{y})$ screened for sign changes on every
accepted integrator step, root location by Brent's method on the step's
dense output to a configurable time tolerance ($\qty{1e-9}{s}$ by default),
direction filters, and integrator restart at the located root with an
optional discrete state update. This is the mechanism by which the mission
sequence acts on the trajectory: the Phase~2 event set is timers and apsis
crossings, and later phases add staging, ignition/cutoff, sphere-of-influence
transitions, node crossings, and ground impact as additional event functions
--- the framework itself does not change. The implementation lives in
\texttt{cpp/src/events.cpp}, on top of the integrators and dense output of
Chapter~\ref{ch:integrators}.

\section{Assumptions and domain of validity}
\label{sec:events:domain}

Assumptions:
\begin{enumerate}
  \item Each event function $g(t, \vv{y})$ is continuous in $t$ along the
    trajectory within a step, and differentiable at its roots with
    $\dot{g}(t^{*}) \neq 0$ (simple crossings). Brent's method still
    converges on roots of odd multiplicity (bisection fallback), but the
    timing-error analysis below assumes simple roots.
  \item Event functions are computable from $(t, \vv{y})$ alone --- no
    history, no side effects. Discrete behavior belongs in the update hook,
    which runs only at located roots.
  \item At most one sign change of each event function occurs per accepted
    step. Sign screening is blind to an even number of crossings inside one
    step (the endpoint signs match); the mitigation is capping the step
    length $h_{\max}$ below half the shortest oscillation period of any
    registered event function.
\end{enumerate}

Domain of validity and out-of-domain behavior, matching the code:
\begin{itemize}
  \item Located event times inherit the dense-output interpolation error
    (equation~\eqref{eq:events:timing}); the framework honors its
    $\qty{1e-9}{s}$ root tolerance with respect to the \emph{interpolated}
    event function. Consumers needing a specific accuracy against the true
    trajectory must budget the $h^{4}$ term via $h_{\max}$, as the apsis
    validation case does.
  \item The apsis event function (equation~\eqref{eq:events:apsis})
    degenerates as $e \to 0$: on a circular orbit $g = \vv{r}\cdot\vv{v}$
    is identically zero and apsides are undefined. The framework itself is
    unaffected (a $g$ that stays at zero never screens as a crossing, by
    the strict-sign rule of equation~\eqref{eq:events:screen}); the event
    simply reports nothing, which is the correct degenerate answer.
  \item An event whose $g$ is exactly zero at the propagation start (or at
    a restart) is not fired until $g$ leaves zero: screening requires a
    strictly signed starting value.
  \item Roots of the same event closer together than the time tolerance are
    indistinguishable by construction and fire once.
\end{itemize}

\section{Derivation}
\label{sec:events:derivation}

\subsection{Sign screening with direction filters}

Let $g_0 = g(t_0, \vv{y}_0)$ and $g_1 = g(t_1, \vv{y}_1)$ be the event
function at the endpoints of an accepted step. A crossing is flagged in
$(t_0, t_1]$ when
\begin{equation}
  \underbrace{g_0 < 0 \;\wedge\; g_1 \ge 0}_{\text{increasing}}
  \qquad\text{or}\qquad
  \underbrace{g_0 > 0 \;\wedge\; g_1 \le 0}_{\text{decreasing}},
  \label{eq:events:screen}
\end{equation}
filtered by the event's direction setting (increasing only, decreasing
only, or either). By the intermediate value theorem the continuous
interpolated event function of
equation~\eqref{eq:events:ghat} then has at least one root in the step.

\subsection{Root location on the dense output}

Evaluating $g$ on the true solution inside a step is not possible without
further integration; the framework instead locates roots of the
\emph{interpolated} event function
\begin{equation}
  \hat{g}(\tau) = g\bigl(\tau,\, \vv{u}(\tau)\bigr),
  \qquad \tau \in [t_0, t_1],
  \label{eq:events:ghat}
\end{equation}
where $\vv{u}$ is the cubic Hermite dense output of
equation~\eqref{eq:integrators:hermite}. Because $\vv{u}$ reproduces the
step endpoints exactly, $\hat{g}$ carries the already-computed endpoint
values $g_0$, $g_1$, and the screened bracket transfers verbatim.

The root is located with Brent's method~\cite{brent1973}: the
guaranteed-convergent combination of bisection, the secant rule, and
inverse quadratic interpolation, taking the interpolation step only when it
falls safely inside the current bracket and shrinks faster than the
previous-but-one step, and bisecting otherwise. Brent proves (i)
convergence is never worse than a bounded multiple of bisection ---
at most $\mathcal{O}\bigl(\log_2^2((b-a)/\delta)\bigr)$ function evaluations
for bracket $[a,b]$ and tolerance $\delta$ --- and (ii) on simple roots of
smooth functions convergence is superlinear~\cite{brent1973}. The
implementation iterates until the bracket width is within
$\delta = \delta_{t} + 4\,\varepsilon\,|t|$, honoring the configured
tolerance $\delta_t$ down to the resolution of double precision at the
root's magnitude.

\subsection{Event-timing error law}
\label{sec:events:timing}

The located time $\hat{t}$ is a root of $\hat{g}$, not of the true
$g^{*}(\tau) = g(\tau, \vv{y}(\tau))$. For a simple root at $t^{*}$,
expanding $g^{*}$ about $t^{*}$ and bounding
$|\hat{g} - g^{*}| \le L_g\,\norm{\vv{u} - \vv{y}}_{\infty}$ with the
Hermite bound of equation~\eqref{eq:integrators:hermite:err} gives
\begin{equation}
  \bigl|\hat{t} - t^{*}\bigr|
    \;\lesssim\;
    \frac{L_g}{|\dot{g}(t^{*})|}\,
    \frac{h^{4}}{384}\, \max \bigl|\vv{y}^{(4)}\bigr|
    \;+\; \delta_t,
  \label{eq:events:timing}
\end{equation}
where $L_g$ is the Lipschitz constant of $g$ in $\vv{y}$ on the step and
$\delta_t$ the root tolerance. Event timing therefore improves as $h^{4}$
and is controlled by $h_{\max}$, exactly as anticipated in
Chapter~\ref{ch:integrators}. For the apsis validation case below
($h_{\max} = \qty{5}{s}$ on the reference orbit), the $h^4$ term evaluates
to $\sim\qty{5e-11}{s}$, so the located times are dominated by the
$\qty{1e-9}{s}$ root tolerance --- three orders inside the
$\qty{1}{\micro s}$ acceptance gate, and consistent with the measured
$3.4\times10^{-4}\,\unit{\micro s}$ worst case.

\subsection{Restart with discrete update}

Reading the event-time state off the dense output would inject the full
$\mathcal{O}(h^4)$ interpolation error into the trajectory. Instead, after
locating $\hat{t}$ the driver \emph{re-integrates} the step from $t_0$ with
step length $\hat{t} - t_0$, so the state delivered at the event carries
integrator accuracy; the shortened step needs no error re-test because it
is interior to a step that already passed. The event's discrete update hook
may then modify the state in place (mass drop, impulsive maneuver, origin
re-centering in later phases), propagation restarts from
$(\hat{t}, \vv{y}(\hat{t}))$, and all event functions are re-evaluated
there. Terminal events stop propagation at $\hat{t}$ instead.

Two safeguards make restart robust at the root. First, the driver fires at
the bracket endpoint on the \emph{post-crossing} side of the root (both
endpoints are within the tolerance), so integration resumes with the sign
change already behind it. Second, for events without a discrete update the
stored screening sign is forced to the crossed side, so interpolant-level
noise around zero cannot re-fire the event; when an update ran, the
re-evaluated value stands, because the hook may legitimately move $g$
(re-arming a timer, changing regime).

\subsection{Phase 2 event functions}

\emph{Timer.} An event at fixed time $t_{\mathrm{e}}$ uses
\begin{equation}
  g(t, \vv{y}) = t - t_{\mathrm{e}},
  \label{eq:events:timer}
\end{equation}
linear in $t$ and independent of the state, so
equation~\eqref{eq:events:timing} reduces to the root tolerance alone
($L_g = 0$): timers isolate the root-finding and restart machinery in
validation.

\emph{Apsis crossing.} For a translational state $[\vv{r}, \vv{v}]$ about a
central body, apsides are the extrema of $r = \norm{\vv{r}}$, i.e. the
zeros of $\dot{r}\,r = \vv{r}\cdot\vv{v}$:
\begin{equation}
  g(t, \vv{y}) = \vv{r} \cdot \vv{v}.
  \label{eq:events:apsis}
\end{equation}
Under two-body dynamics (equation~\eqref{eq:twobody:accel}),
$\dot{g} = \vv{v}\cdot\vv{v} + \vv{r}\cdot\dot{\vv{v}} = v^2 - \mu/r$.
Evaluating on the ellipse with semi-latus rectum $p = a(1 - e^2)$
(Vallado~\cite{vallado2013}): at periapsis $v_p^2 = (\mu/p)(1+e)^2$ and
$\mu/r_p = (\mu/p)(1+e)$, so $\dot{g}_p = \mu e (1 + e)/p > 0$; at apoapsis
$\dot{g}_a = -\mu e (1 - e)/p < 0$. Periapsis is therefore an increasing
crossing of $g$ and apoapsis a decreasing one, which is how the two
direction-filtered event registrations in the validation suite distinguish
them, and $|\dot g|$ at the roots --- $\sim\qty{8.8e6}{m^2/s^3}$ for the
reference orbit --- is the denominator in
equation~\eqref{eq:events:timing}.

\section{Discretization and implementation}
\label{sec:events:implementation}

\begin{enumerate}
  \item \textbf{Module and traceability.} The framework lives in
    \texttt{cpp/src/events.cpp} with the public interface in
    \texttt{cpp/include/star/events.hpp}. Source comments echo the labels
    \texttt{eq:events:screen} and \texttt{eq:events:ghat} at the screening
    and root-location code, and \texttt{eq:events:timer} and
    \texttt{eq:events:apsis} at the two event functions.
  \item \textbf{Extension contract.} An event type is a callable
    $g(t, \vv{y})$ plus a direction, an optional update hook, and a
    terminal flag (\texttt{EventSpec}); Phase 3+ event types plug in
    without framework changes. All framework references to user code are
    non-owning function references, matching the integrator's hot-loop
    interface.
  \item \textbf{Determinism (D-10).} Events are screened in registration
    order every step; when several events root within the same step the
    earliest located time fires, with exact ties resolved to the lowest
    index. Brent's iteration is pure double arithmetic on the bracket; all
    workspace is pre-allocated. The only allocation after setup is the
    growth of the caller-supplied event log when an event actually fires
    (callers may reserve).
  \item \textbf{Bracket endpoints for free.} Because the dense output
    reproduces step endpoints bitwise, the screening values $g_0$, $g_1$
    seed Brent's bracket without re-evaluation, and $g_1$ of an accepted
    step is reused as $g_0$ of the next.
  \item \textbf{Tolerance floor.} The root tolerance degrades gracefully to
    $4\varepsilon|t|$ when the configured $\delta_t$ is below the
    resolution of double at the root's magnitude.
\end{enumerate}

\section{Validation evidence}
\label{sec:events:validation}

Table~\ref{tab:events:validation} lists the tests that constitute this
chapter's validation evidence. Each identifier is machine-checked to exist
in the test tree by \texttt{scripts/lint\_chapters.py}; the analytic apsis
times are committed goldens with provenance
(\texttt{tests/golden/integrators/manifest.toml}), and the apsis
measurement is reproduced from Python by
\texttt{tests/python/test\_integrators.py}.

\begin{table}[htbp]
  \centering
  \caption{Validation evidence for the event-detection framework. Measured
    values are from the acceptance run recorded in the test sources
    (MSVC x64, \texttt{/fp:strict}).}
  \label{tab:events:validation}
  \begin{tabular}{lp{8.6cm}}
    \toprule
    Test ID & Acceptance basis \\
    \midrule
    \testid{brent_root_known_functions} &
      Brent's method against known roots (quadratic, transcendental,
      triple root, exact-zero endpoint, non-bracketing rejection): located
      within $\delta_t + 4\varepsilon|t|$, iteration counts consistent with
      superlinear convergence~\cite{brent1973}. \\
    \testid{timer_event_location} &
      Timer events (equation~\eqref{eq:events:timer}) located within the
      $\qty{1e-9}{s}$ root tolerance exactly; terminal event stops
      propagation at the located root. \\
    \testid{event_restart_discrete_update} &
      An impulsive along-track $\Delta v$ applied by an update hook at a
      timer event: final state matches the two-segment analytic Kepler
      composition to $\qty{2e-3}{m}$ / $\qty{2e-6}{m/s}$; the event fires
      exactly once (no duplicate re-fire at the root). \\
    \testid{apsis_event_times_analytic} &
      Periapsis and apoapsis times (equation~\eqref{eq:events:apsis},
      direction-filtered) over 3.2 orbits of the eccentric reference orbit
      vs the committed analytic times: worst error
      $3.4\times10^{-4}\,\unit{\micro s}$, gate $< \qty{1}{\micro s}$
      (Phase 2 exit criterion 5). \\
    \bottomrule
  \end{tabular}
\end{table}

\section{References}
\label{sec:events:references}

This chapter relies on Brent~\cite{brent1973} for the root-finding
algorithm and its convergence guarantees, on Hairer, N{\o}rsett, and
Wanner~\cite{hairer1993} for the dense output the roots are located on, and
on Vallado~\cite{vallado2013} for the apsis geometry of the ellipse. Full
bibliographic details appear in the bibliography.
